\section{An Explicit Dyadic Bound for the Retained Class-Mean Method}\label{sec:dyadic52}
The resource-path guarantee does not require class-mean subdivision. We nevertheless give an explicit bound for that retained method, both for comparison and to remove ambiguity about its stopping scale.

Let $X=[\ell,u]$ be a projected class-mean box intersecting $\sum_gW_gs_g=B$. Fix a dependent class $g_0$ and let $d=E-1>0$. A zero-price box oracle supplies an upper value. Its selected book is feasible at the original promise; bounded one-direction transfers of its class means restore equality. Since every fixed-book class response is $L$-Lipschitz,
\[
 \text{zero-price upper value}-\text{recovered lower value}
 \leq L\sum_gW_g(u_g-\ell_g).
\]
Projection onto the root equality gives
\[
 W_{g_0}(u_{g_0}-\ell_{g_0})\leq\sum_{g\ne g_0}W_g(u_g-\ell_g).
\]
Thus the explicit free-coordinate stopping threshold
\begin{equation}\label{eq:etabox52}
 \eta_{\rm box}=\frac{\varepsilon}{2L(E-1)}
\end{equation}
suffices: if every free weighted width is at most $\eta_{\rm box}$, the gap is at most $\varepsilon$. The zero price must be included even when additional prices are used.

\begin{proposition}[Dyadic cover, with the ceiling retained]\label{prop:dyadic52}
Let $w_g^0=W_g(u_g^0-\ell_g^0)$ be the initial free weighted widths. Bisect a largest free weighted width at every unresolved node, projecting children onto the equality. Define
\[
 q_g=\begin{cases}
 0,&w_g^0=0,\\
 \max\{0,\lceil\log_2(w_g^0/\eta_{\rm box})\rceil\},&w_g^0>0.
 \end{cases}
\]
A complete binary cover has at most
\begin{equation}\label{eq:dyadic52}
 M\leq\prod_{g\ne g_0}2^{q_g}
 \leq\left(\max\left\{1,\frac{4L(E-1)}{\varepsilon}\right\}\right)^{E-1}
\end{equation}
leaves, and at most $2M-1$ nodes when both children of every split are counted. The first product is a dyadic upper bound, not a claim that projection must generate every tensor-grid leaf.
\end{proposition}
\begin{proof}
A node with gap above $\varepsilon$ has some free width greater than $\eta_{\rm box}$. Its largest free width also exceeds this threshold. A split halves that coordinate's width, while projection never increases any width. Coordinate $g$ can therefore be split at most $q_g$ times on a root-to-leaf path. Every path has depth at most $\sum_gq_g$, so its binary tree has at most $2^{\sum_gq_g}$ leaves. A full binary tree with $M$ leaves has $2M-1$ nodes. If infeasible children are pruned, including them before pruning provides the same upper bound.

For the coarse bound, each $w_g^0\leq1$. If $\eta_{\rm box}\geq1$, all $q_g$ vanish. Otherwise, $2^{\lceil\log_2(1/\eta_{\rm box})\rceil}\leq2/\eta_{\rm box}=4L(E-1)/\varepsilon$. Multiplication proves \eqref{eq:dyadic52}.
\end{proof}

For width one and threshold $1/3$, the dyadic requirement is four leaves, not three. The coarse bound must be evaluated with the \emph{same} threshold: \eqref{eq:etabox52} then gives $4L(E-1)/\varepsilon=6$ in one free dimension. Conversely, setting $4L/\varepsilon=3$ gives $\eta_{\rm box}=2/3$ and requires only two leaves. These examples distinguish an omitted ceiling in a putative exact product from a valid coarser bound that already includes its factor of two. They also explain why the stopping threshold must not silently change from $\varepsilon/(2Ld)$ to $\varepsilon/(4Ld)$.

A repaired class-mean net remains valid as well. For each feasible anchor, round $E-1$ weighted class means down with step $\varepsilon/(2L(E-1))$, set the final coordinate by the equality, clip and repair the remaining deficit by bounded transfers. The weighted absolute perturbation is at most twice the rounded mass. Exact conditional class optimization then loses at most $\varepsilon$. This yields at most $N(1+2L(E-1)/\varepsilon)^{E-1}$ calls. Its class-dependent exponent is a property of this particular enumeration, not a lower bound on the joint problem; Theorem~\ref{thm:approx52} removes that exponent.
